%!TEX root = ExerciseSheet1.tex
% -----------------------------------------------------------------------------
% PART A
% -----------------------------------------------------------------------------{
\subsection*{Part A: The Idea}
\textit{These questions test the mathematical definitions and calculations devoid of physical context.}

\begin{enumerate}[label=\textbf{Q\arabic*.}, leftmargin=*]
    \item \textbf{Data Types} \\
    Classify the following abstract variables as Qualitative, Quantitative Discrete, or Quantitative Continuous:
    \begin{enumerate}[label=\arabic*.]
        \item The number of elements in a set.
        \item A label assigned to a category (e.g., \lq Pass\rq , \lq Fail\rq ).
        \item A real number generated by a continuous mathematical function.
    \end{enumerate}
    \vspace{0.3cm}

    \item \textbf{Measures of Central Tendency} \\
    Consider the following abstract dataset of values: 
    \[ \{ 2, 3, 4, 4, 4, 4 ,5, 10, 52 \} \]
    \begin{enumerate}[label=\arabic*.]
        \item Calculate the sample mean ($\bar{x}$).
        \item Determine the median and the mode.
        \item If the value $52$ was a data entry error and should have been $12$, calculate the new mean and median. Which measure was more affected by the outlier?
    \end{enumerate}
    \vspace{0.3cm}

    \item \textbf{Measures of Dispersion} \\
    Consider a small sample dataset: 
    \[ \{ 2, 4, 6, 8, 10 \} \]
    \begin{enumerate}[label=\arabic*.]
        \item Calculate the range of the data.
        \item Calculate the sample variance ($s^2$). Be sure to explicitly show the use of Bessel's Correction.
        \item Calculate the sample standard deviation ($s$). 
        \item If this data represented an entire population ($N=5$) rather than a sample, what would the population variance ($\sigma^2$) be?
    \end{enumerate}

\end{enumerate}

\vspace{0.8cm}

% -----------------------------------------------------------------------------
% PART B
% -----------------------------------------------------------------------------
\subsection*{Part B: Exam style questions: Engineering Application}
\textit{These questions test the exact same mathematical concepts as Part A, but applied to the Electrical Engineering scenarios discussed in the lecture.}

\begin{enumerate}[label=\textbf{Q\arabic*.}, leftmargin=*, start=4]

    \item \textbf{Data Types in the Lab} \\
    Classify the following engineering measurements as Qualitative, Quantitative Discrete, or Quantitative Continuous:
    \begin{enumerate}[label=\arabic*.]
        \item The number of defective pixels on an LCD monitor.
        \item The exact current (in mA) drawn by a microcontroller during sleep mode.
        \item The colour band on a resistor (e.g., Red, Black, Orange).
    \end{enumerate}
    \vspace{0.3cm}

    \item \textbf{Central Tendency \& Network Latency} \\
    An engineer measures the round-trip latency (in ms) of $8$ data packets sent to a remote server. The results are: 
    \[ \{ 4\text{ ms}, 5\text{ ms}, 6\text{ ms}, 6\text{ ms}, 7\text{ ms}, 9\text{ ms}, 10\text{ ms}, 42\text{ ms} \} \]
    \begin{enumerate}[label=\arabic*.]
        \item Calculate the mean, median, and mode latency.
        \item The $42\text{ ms}$ reading was caused by a temporary network routing spike. As an engineer summarising the ``typical'' performance of this network, would you report the mean or the median? Justify your choice based on the data distribution.
    \end{enumerate}
    \vspace{0.3cm}

    \item \textbf{Component Tolerance \& Spread} \\
    You pull $5$ nominally identical $6\,\Omega$ resistors from a bin and measure their exact resistance. The readings are: 
    \[ \{ 3\,\Omega, 4\,\Omega, 4\,\Omega, 5\,\Omega, 6\,\Omega \} \]
    \textit{(Note: These are terrible resistors!)}
    \begin{enumerate}[label=\arabic*.]
        \item Calculate the sample variance of this batch. What are the physical units of your answer?
        \item Calculate the sample standard deviation. What are the physical units of this answer?
        \item Referring to your answers regarding units, explain why engineers prefer to communicate uncertainty using Standard Deviation rather than Variance.
        \item On the other hand mathematician prefer Variance, why?
    \end{enumerate}
    \vspace{0.3cm}

    \item \textbf{Sampling \& Bias} \\
    A factory produces 10,000 lithium-ion batteries per day. To test the reliability of the daily batch, the Quality Assurance team selects the first $50$ batteries produced right after the machines are turned on at 8:00 AM. They find $0$ defects and conclude the entire daily production is 100\% reliable.
    \begin{enumerate}[label=\arabic*.]
        \item Is this a representative sample? Why or why not?
        \item Identify a specific variable that might introduce bias into this sampling method. \textit{(Hint: Think about machine temperature or calibration over a 24-hour shift)}.
    \end{enumerate}

\end{enumerate}




